%! Unit 7: The Singular Value Decomposition — Summative Assessment — Unit Test (blank)
\documentclass[10pt]{article}
\usepackage{linalg-article}
\usepackage{linalg-boxes}
\newcommand{\xx}{\mathbf{x}}
\newcommand{\vv}{\mathbf{v}}
\newcommand{\uu}{\mathbf{u}}
\newcommand{\T}{^{\mathsf T}}

% Part-divider strip
\newcommand{\parthead}[1]{%
  \vspace{6pt}\noindent
  \begin{headlinebox}{burgundy}{\color{white}\bfseries #1}\end{headlinebox}%
  \vspace{4pt}}

\begin{document}

\pageheader{Unit 7: The Singular Value Decomposition}{Unit Test}
\namedateperiod

\vspace{0.06in}
\noindent\textbf{Instructions:} Show all work. No notes unless your teacher says so.

% ======================================================================
\parthead{Part A --- Vocabulary (8 pts)}
Write the letter of the best description next to each term.

\vspace{4pt}
\noindent
\begin{minipage}[t]{0.40\textwidth}
\begin{enumerate}[leftmargin=2.2em, itemsep=7pt, label=\arabic*.]
  \item Singular Value Decomposition \blank{1.2cm}
  \item Orthogonal matrix \blank{1.2cm}
  \item Singular value \blank{1.2cm}
  \item Energy of a matrix \blank{1.2cm}
  \item Outer product \blank{1.2cm}
  \item Principal component \blank{1.2cm}
  \item Singular vectors \blank{1.2cm}
  \item Rank-one matrix \blank{1.2cm}
\end{enumerate}
\end{minipage}%
\hfill
\begin{minipage}[t]{0.57\textwidth}
\small
\begin{description}[leftmargin=1.4em, itemsep=3pt]
  \item[(A)] The factorization $A=U\Sigma V\T$ (rotate--stretch--rotate) that \emph{every} matrix has.
  \item[(B)] A square matrix with perpendicular unit columns, so $Q\T Q=I$ and $Q^{-1}=Q\T$.
  \item[(C)] A number $\sigma\ge0$: how far $A$ stretches along a singular direction (a half-width of the ellipse).
  \item[(D)] The total $\sigma_1^2+\sigma_2^2+\cdots$, equal to the sum of the squares of all entries of $A$.
  \item[(E)] A column times a row, $\uu\vv\T$ --- a whole matrix built from two vectors.
  \item[(F)] A top singular direction of centered data --- the direction of greatest spread.
  \item[(G)] The perpendicular unit inputs $\vv_i$ and outputs $\uu_i$ with $A\vv_i=\sigma_i\uu_i$.
  \item[(H)] A matrix whose every column is a multiple of one vector (like $\uu\vv\T$).
\end{description}
\end{minipage}

% ======================================================================
\parthead{Part B --- Multiple Choice (12 pts)}
Circle the best answer.

\begin{enumerate}[leftmargin=1.9em, itemsep=9pt]
  \item Which matrices have a Singular Value Decomposition?
    \hfill (A) every matrix (any shape) \quad (B) only symmetric ones \quad (C) only square ones \quad (D) only invertible ones
  \item The input singular vectors $\vv_i$ are eigenvectors of
    \hfill (A) $A$ \quad (B) $AA\T$ \quad (C) $A\T A$ \quad (D) $A^{-1}$
  \item Once you have $\sigma_i$ and $\vv_i$, the output vector is
    \hfill (A) $\uu_i=A\vv_i/\sigma_i$ \quad (B) $\uu_i=\sigma_i A\vv_i$ \quad (C) $\uu_i=A\T\vv_i$ \quad (D) $\uu_i=\vv_i/\sigma_i$
  \item The outer product $\uu\vv\T$ is a
    \hfill (A) scalar \quad (B) rank-one matrix \quad (C) diagonal matrix \quad (D) inverse
  \item An orthogonal matrix preserves
    \hfill (A) length: $|Q\xx|=|\xx|$ \quad (B) always $\det=+1$ \quad (C) nothing \quad (D) only direction
  \item In PCA, the first principal component is the direction of
    \hfill (A) least spread \quad (B) greatest spread \quad (C) zero spread \quad (D) the mean
\end{enumerate}

% ======================================================================
\parthead{Part C --- Short Answer \& Computation (35 pts)}
Show your work.

\begin{enumerate}[leftmargin=1.9em, itemsep=12pt]
  \item \textbf{(a)} The diagonal matrix $D=\begin{bsmallmatrix}5&0\\0&2\end{bsmallmatrix}$ acts on the unit circle.
        What are its singular values $\sigma_1\ge\sigma_2$, and what are the half-widths of the ellipse it makes?
        \textbf{(b)} Find the singular values of the swap $S=\begin{bsmallmatrix}0&3\\2&0\end{bsmallmatrix}$.
        \vspace{2.4cm}
  \item Find the \textbf{singular values} of $A=\begin{bsmallmatrix}3&1\\-1&-3\end{bsmallmatrix}$ by forming
        $A\T A$, finding its eigenvalues, and taking $\sigma=\sqrt{\lambda}$.
        \vspace{2.8cm}
  \item For that same $A=\begin{bsmallmatrix}3&1\\-1&-3\end{bsmallmatrix}$, the input vector for $\sigma_1=4$ is
        $\vv_1=\tfrac{1}{\sqrt2}\begin{bsmallmatrix}1\\1\end{bsmallmatrix}$. Find the \textbf{output singular vector}
        $\uu_1=\dfrac{A\vv_1}{\sigma_1}$, and check that it is a unit vector.
        \vspace{2.8cm}
  \item \textbf{(a)} Compute the outer product $\uu\vv\T$ where $\uu=\begin{bsmallmatrix}2\\1\end{bsmallmatrix}$ and
        $\vv=\begin{bsmallmatrix}1\\3\end{bsmallmatrix}$. \textbf{(b)} Explain why the result is a \textbf{rank-one} matrix.
        \vspace{2.6cm}
  \item The symmetric matrix $A=\begin{bsmallmatrix}6&4\\4&6\end{bsmallmatrix}$ has singular values $\sigma_1=10,\ \sigma_2=2$
        with $\uu_1=\vv_1=\tfrac{1}{\sqrt2}\begin{bsmallmatrix}1\\1\end{bsmallmatrix}$.
        \textbf{(a)} Build the best rank-one approximation $A_1=\sigma_1\uu_1\vv_1\T$.
        \textbf{(b)} What fraction of the \textbf{energy} does $A_1$ keep?
        \vspace{2.8cm}
  \item Let $Q=\dfrac15\begin{bsmallmatrix}4&-3\\3&4\end{bsmallmatrix}$.
        \textbf{(a)} Show $Q$ is \textbf{orthogonal} by checking $Q\T Q=I$.
        \textbf{(b)} Write down $Q^{-1}$.
        \textbf{(c)} The vector $\xx=\begin{bsmallmatrix}0\\5\end{bsmallmatrix}$ has length $5$. Find $Q\xx$ and its length, and explain.
        \vspace{2.8cm}
  \item Four data points are $(6,5),\ (5,6),\ (2,3),\ (3,2)$.
        \textbf{(a)} \textbf{Center} the data (subtract the mean). \textbf{(b)} Form $A\T A$.
        \textbf{(c)} Find its eigenvalues and the first principal component $\mathrm{PC}_1$.
        \textbf{(d)} What fraction of the spread does $\mathrm{PC}_1$ capture?
        \vspace{3.0cm}
\end{enumerate}

% ======================================================================
\parthead{Part D --- Extended Response (10 pts)}
Explain your reasoning in full sentences.

\begin{enumerate}[leftmargin=1.9em, itemsep=16pt]
  \item \emph{Every} matrix has an SVD, even though not every matrix has real eigenvectors.
        \textbf{(a)} Explain why $A\T A$ is symmetric for \emph{any} matrix $A$, and why that guarantees its
        eigenvalues are real with $\lambda\ge0$ and its eigenvectors are perpendicular.
        \textbf{(b)} Explain why this makes $\sigma=\sqrt{\lambda}$ a real number $\ge0$, so every matrix has singular values.
        \vspace{3.0cm}
  \item The SVD $A=U\Sigma V\T$ is ``rotate--stretch--rotate.''
        \textbf{(a)} Explain what each of the three factors does geometrically, and why \emph{all} the stretching lives in $\Sigma$.
        \textbf{(b)} An orthogonal matrix $Q$ has $Q\T Q=I$. Explain why this makes $Q^{-1}=Q\T$ (a free inverse) and why $Q$ preserves length.
        \vspace{2.8cm}
\end{enumerate}

\end{document}
